% paper/sections/introduction.tex

\section{Introduction}
\label{sec:intro}

A \PID{} controller carefully tuned for nominal \AUV{} physics achieves
$100\%$ goal-reaching success in simulation---and $3\%$ when the same
fluid parameters shift by a modest, ocean-realistic amount.
This single result, demonstrated quantitatively in
Section~\ref{sec:pid_results}, encapsulates the central challenge of
autonomous underwater vehicle (\AUV{}) deployment: controllers
optimised for nominal conditions fail when drag, buoyancy, and water
current vary with biofouling, salinity gradients, and tidal
patterns---as they invariably do in real ocean operations.

Reinforcement learning (\RL{}) trained with domain randomisation (\DR{})
offers a principled solution: by varying physics parameters across
training episodes, policies can generalise to the distribution of
conditions encountered in deployment~\cite{tobin2017,peng2018}.
Several recent works have applied \SAC{} and \DR{} to \AUV{}
control~\cite{arndt2024,chu2025,chaffre2025}, demonstrating
sim-to-real transfer for station-keeping and path-following.
However, standard Uniform \DR{}---sampling uniformly from fixed ranges
throughout training---exposes the agent simultaneously to easy and hard
conditions from the very first episode.
We show this produces policies with high inter-seed variance
($47.67\% \pm 33.83\%$ without any \DR{}),
making deployment reliability unpredictable.

A natural improvement is to schedule \DR{} difficulty: begin with
narrow physics ranges and expand them as the agent demonstrates
competence.
OpenAI's Automatic Domain Randomisation (\textsc{Adr})~\cite{akkaya2019}
applied this principle to dexterous manipulation.
However, no prior work has applied curriculum \DR{} to underwater fluid
physics---a qualitatively distinct parameter space where drag, current,
and buoyancy are nonlinearly coupled---and no prior work evaluates
energy efficiency, a metric of direct operational significance for
battery-constrained \AUV{}s.

We present \emph{Curriculum Domain Randomisation} (\CDR{}) for \AUV{}
fluid physics: a performance-gated strategy that starts from narrow
physics ranges and expands them when a rolling success-rate window
exceeds a threshold, inspired by \textsc{Adr}~\cite{akkaya2019} but
simplified for underwater physics---expanding all six parameters
simultaneously with a single shared window, without per-parameter
buffers or Bayesian models.
We compare \CDR{} against Naive \SAC{}, Uniform \DR{}, and a \PID{}
baseline across nine experiments ($3$ conditions~$\times$~$3$
seeds~$\times$~$10^6$ steps) and evaluate zero-shot transfer to
a held-out test distribution with elevated drag and strong currents.

Our contributions are:
\begin{enumerate}
  \item The \textbf{first systematic comparison} of curriculum versus
    uniform \DR{} strategies for \AUV{} fluid physics, demonstrating
    that \CDR{} produces $1.13\%$ more energy-efficient
    policies at comparable transfer-success rates
    ($98.67\% \pm 1.89\%$).
  \item A \textbf{quantitative \PID{} fragility benchmark} showing
    $100\% \to 3\%$ success collapse under realistic fluid-parameter
    shift, establishing the practical necessity of physics-aware
    training for real \AUV{} deployment.
  \item An \textbf{open-source \MuJoCo{} \AUV{} environment}
    (\textit{Halcyon~X4}) with nine randomisable parameters, released
    for reproducible underwater robotics research at
    \url{https://github.com/Itzlimon22/rl_robotics}.
\end{enumerate}